\chapter{绪论}

\section{研究背景与意义}

随着云计算与微服务架构的广泛应用，现代软件系统在规模、复杂度和运行动态性方面持续提升，其运行、维护和保障（简称"运维"）面临巨大挑战。在故障定位、异常关联和跨源信息整合方面，传统方法已难以满足高效故障诊断需求。据统计，一次复杂故障的平均诊断时间可达数小时，其中大量时间消耗在信息检索与初步排查上。

传统的软件运维主要依赖人力和固定规则脚本，已难以应对动态、分布式的系统环境。传统智能运维（AIOps）方法在处理非结构化日志、理解复杂故障逻辑、进行根本原因推断等方面仍存在不足，具体表现为：（1）传统规则与脚本方法难以处理非结构化日志和复杂故障传播链；（2）传统机器学习方法对领域迁移和跨任务泛化支持不足；（3）在真实运维场景中，根因分析仍高度依赖运维人员经验。

近年来，以GPT、Llama等为代表的大语言模型展现了强大的自然语言理解和知识推理能力，为软件运维中的日志理解、知识检索增强与故障推理提供了新的技术可能。然而，将大语言模型直接应用于专业运维场景仍面临以下核心挑战：

（1）\textbf{领域知识不足问题}。通用大语言模型对特定系统的架构背景、历史故障案例和运维经验了解有限，仅依赖模型参数知识容易出现回答不准确的问题。

（2）\textbf{异构运维数据的上下文组织问题}。日志、监控指标、告警等数据格式各异，如何将异构数据统一组织为适合大模型分析的上下文，是提升分析质量的关键。

（3）\textbf{故障推理稳定性与可解释性问题}。根因分析是一个高度依赖领域知识和推理能力的任务，直接调用大模型输出的结论往往缺乏依据支撑，需要结合外部知识与实时数据进行约束。

针对上述挑战，本文研究一种结合检索增强、多源证据融合与结构化推理的软件运维故障诊断方法。主要研究工作包括：

第一，\textbf{研究面向运维文档的知识组织与检索方法}。通过收集整理运维文档和历史故障案例，提出基于文档结构的文本切分策略，提升领域知识对故障诊断的支撑能力。

第二，\textbf{研究日志与指标的证据融合机制}。通过设计日志与监控指标的关联机制，为大模型构建可推理的结构化上下文。

第三，\textbf{研究基于结构化Prompt的根因分析方法}。通过结构化提示模板引导模型进行逐步推理，以提升输出的一致性、可解释性与可验证性。

第四，\textbf{通过原型系统与对照实验验证上述方法的有效性}。构建端到端的故障诊断流程，验证所提方法在提升故障诊断准确率方面的作用。

\section{国内外研究现状}

当前基于大模型的软件运维研究主要集中在日志分析、检索增强生成和根因分析三个方向。

\subsection{基于大模型的日志分析研究}

日志分析是将大语言模型应用于软件运维的重要切入点。Liu等人提出的LogPrompt方法~\cite{liu2023logprompt}通过设计面向日志分析的提示模板，使大语言模型能够识别异常日志、分类日志类型并生成分析结果，在公开数据集上验证了方法的有效性。Ji等人提出的SuperLog方法~\cite{ji2024superlog}探索了将领域知识注入日志分析流程的方法，通过持续预训练使模型更好地理解运维术语和故障模式。Jiang等人提出的LILAC方法~\cite{jiang2024lilac}通过引入缓存机制提升了日志解析效率。然而，现有日志研究多聚焦单任务分析，难以支撑完整根因分析链路。

\subsection{基于检索增强生成的运维知识增强研究}

检索增强生成（RAG）技术为解决大语言模型领域知识不足问题提供了有效途径。Zhang等人提出的RAG4ITOps框架~\cite{zhang2024rag4itops}构建了面向IT运维的知识库，将检索增强生成技术应用于运维问答任务，在企业数据集上验证了方法的有效性。Jiang等人提出的Active RAG方法~\cite{jiang2023active}探索了动态检索策略，增强了知识获取的针对性。Isaza等人~\cite{isaza2024ragincident}将RAG技术应用于IT支持工单推荐，为知识库构建提供了参考。然而，现有RAG研究多关注问答与推荐，缺少面向运维故障诊断的完整多源推理验证。

\subsection{基于大语言模型的根因分析研究}

根因分析（Root Cause Analysis，RCA）是将大语言模型应用于运维诊断的核心任务。Roy等人~\cite{roy2024agents}探索了基于大语言模型代理的根因分析方法，验证了大模型在多步推理方面的潜力。Szandała等人~\cite{szandala2025reliability}评估了大语言模型在混沌工程场景下的故障诊断能力，指出推理过程的透明性和可解释性是当前方法的主要瓶颈。Pham等人~\cite{pham2024rcacausal}对基于因果推断的微服务根因分析方法进行了系统性评测，指出传统方法在复杂故障传播场景下的局限性。然而，现有RCA研究常依赖单一数据视角，缺少将知识检索、实时数据和结构化推理相结合的完整方法验证。

\subsection{研究缺口总结}

综合上述研究现状可以看出，现有工作主要存在以下不足：（1）现有日志研究多聚焦单任务分析，难以支撑完整根因分析链路；（2）现有RCA研究常依赖单一数据视角，缺少对多源证据的系统性融合；（3）现有RAG研究多关注问答与推荐，缺少面向运维故障诊断的完整方法链路验证。本文正是围绕这一研究缺口展开，探索一种融合领域知识检索、实时数据融合与结构化提示引导的故障诊断方法。

\section{研究内容}

围绕上述研究背景与现状分析，本文重点开展以下研究工作：

第一，研究面向运维领域的知识库构建方法。通过收集整理运维文档（如标准操作流程、历史故障案例等），提出基于文档结构的文本切分策略，使用向量数据库实现运维知识的存储与检索。

第二，研究证据融合方法。通过设计日志与监控指标的关联机制，将异构数据统一组织为结构化的分析上下文，为后续推理提供依据。

第三，研究基于结构化Prompt的根因分析推理方法。设计结构化提示模板，引导大模型进行逐步推理，提升分析结果的可解释性和准确性。

第四，构建原型系统并基于典型故障场景开展对照实验，用于验证所提方法的有效性。

\section{研究方法与技术路线}

本文采用理论分析与原型验证相结合的研究方法，主要包括以下方面。

在理论层面，首先通过文献调研系统梳理大语言模型、检索增强生成和根因分析等领域的研究进展，提炼出现有方法的不足与本文方法的切入点；其次采用方法设计法，针对领域知识不足、异构数据难组织和推理结果不稳定三个问题分别设计相应方法；最后采用对比实验法，通过消融分析量化各模块对诊断效果的贡献。

在技术层面，采用Python语言开发原型系统。知识库基于ChromaDB向量数据库构建，向量嵌入使用预训练语言模型。大语言模型调用采用OpenAI API接口。原型系统前端基于Streamlit框架，后端基于FastAPI框架。

本文技术路线如下：首先进行需求分析与知识库构建；其次完成多源数据处理与融合；然后进行根因分析引擎设计；最后通过测试案例验证系统效果。

本文的技术路线图如图~\ref{fig:roadmap} 所示。

\begin{figure}[htbp]
\centering
\begin{tikzpicture}[
    >=stealth,
    node distance=1.5cm,
    font=\songti\zihao{-5},
    proc/.style={draw, rounded corners, fill=blue!10, minimum width=3cm, minimum height=0.8cm, align=center},
    arrow/.style={->, thick}
]

\node[proc] (step1) at (0, 3) {问题分析};
\node[proc] (step2) at (0, 1.5) {相关工作梳理};
\node[proc] (step3) at (0, 0) {方法设计};
\node[proc] (step4) at (0, -1.5) {原型实现};
\node[proc] (step5) at (0, -3) {对比实验};
\node[proc] (step6) at (0, -4.5) {结果分析};

\draw[arrow] (step1) -- (step2);
\draw[arrow] (step2) -- (step3);
\draw[arrow] (step3) -- (step4);
\draw[arrow] (step4) -- (step5);
\draw[arrow] (step5) -- (step6);

\end{tikzpicture}
\caption{本文技术路线图}
\label{fig:roadmap}
\end{figure}

\section{论文结构安排}

本文共分为五章，各章内容安排如下：

第一章为绪论，介绍课题的研究背景与意义，综述国内外研究现状并总结研究缺口，明确本文研究内容与方法，概括论文结构安排。

第二章为相关技术基础，介绍智能运维发展脉络、大语言模型、检索增强生成、多源运维数据分析和根因分析的核心概念与技术原理，为后续方法设计奠定基础。

第三章为方法设计与原型实现，围绕故障诊断中知识不足、数据异构和推理不稳定三个问题，从知识增强、多源上下文构建和结构化推理三个方面详细说明本文方法设计与原型实现。

第四章为实验与结果分析，通过典型故障场景测试验证所提方法的效果，并从多维度分析实验结果，总结系统的不足与改进方向。

第五章为总结与展望，总结本文的研究工作与主要创新点，展望未来的研究方向。
